\chapter{Contexte g\'en\'eral du projet}
\label{chap:contexte}

\section{Introduction}
% Courte intro (3-4 lignes) annoncant le contenu du chapitre :
% organisme d'accueil, etude de l'existant, problematique, solution, methodologie.

[\`A compl\'eter]

\section{Pr\'esentation de l'organisme d'accueil}
% A completer avec les informations reelles de l'entreprise :
%   - Nom, secteur d'activite, date de creation
%   - Domaines d'expertise / clients
%   - Organigramme si pertinent
%   - Lien avec le projet (pourquoi cette entreprise developpe E-FORCE)

[\`A compl\'eter : pr\'esentation de l'entreprise / organisme d'accueil --
historique, activit\'es, organisation.]

\section{Pr\'esentation g\'en\'erale du projet}
% Vue d'ensemble du projet E-FORCE (peut reprendre/adapter le contenu
% de DOCUMENTATION.md) :
%   - Plateforme de gestion du commerce international
%   - 5 grands modules : authentification/Keycloak, processus BPMN,
%     regles metier + DRL, module IA/NLP, import/export de dossiers
%   - Acteurs : SuperAdmin, Gestionnaire des processus metier,
%     Gestionnaire des regles metier

La plateforme \textbf{E-FORCE} est une application web destin\'ee \`a la
gestion du commerce international et des proc\'edures douani\`eres. Elle
s'organise autour de cinq grands modules~:

\begin{itemize}
    \item \textbf{Authentification et gestion des utilisateurs} : bas\'ee
    sur Keycloak (SSO), avec gestion des r\^oles (SuperAdmin,
    Gestionnaire des processus m\'etier, Gestionnaire des r\`egles
    m\'etier) ;
    \item \textbf{Gestion des processus m\'etier (BPMN)} : cr\'eation,
    modification et visualisation de processus repr\'esent\'es sous forme
    de diagrammes BPMN 2.0 ;
    \item \textbf{Moteur de r\`egles m\'etier} : d\'efinition de r\`egles
    (conditions + actions) avec g\'en\'eration automatique de code
    \textit{Drools Rule Language} (DRL) et gestion du versioning ;
    \item \textbf{Module IA / NLP} : conversion automatique d'une phrase
    en langage naturel en r\`egle m\'etier structur\'ee et en code DRL,
    avec calcul d'un score de confiance ;
    \item \textbf{Gestion des dossiers d'import/export} : suivi des
    dossiers de marchandises import\'ees/export\'ees, avec export des
    donn\'ees au format PDF et Excel.
\end{itemize}

\section{\'Etude de l'existant}
% Decrire comment les taches sont realisees AVANT ce projet :
%   - outils utilises actuellement (Excel, papier, code DRL ecrit a la main, etc.)
%   - acteurs impliques et leurs difficultes quotidiennes

[\`A compl\'eter : d\'ecrire la situation actuelle / les solutions existantes
et leurs limites.]

\subsection{Critique de l'existant}
% Lister les limites/problemes identifies, par exemple :
%   - Necessite de competences techniques (DRL) pour creer une regle
%   - Absence de visualisation des processus metier
%   - Gestion manuelle / dispersee des dossiers import-export
%   - Pas de tracabilite des modifications de regles (historique)

[\`A compl\'eter : limites de la solution actuelle.]

\subsection{Solution propos\'ee}
% Expliquer en quoi E-FORCE repond a chaque limite identifiee ci-dessus
% (tableau "probleme -> solution" recommande)

[\`A compl\'eter : pr\'esenter la solution propos\'ee en r\'eponse \`a la
critique de l'existant -- \'eventuellement sous forme de tableau
probl\`eme/solution.]

\section{M\'ethodologie de gestion de projet}
% Decrire la methodologie suivie (Scrum, Kanban, cycle en V, etc.) :
%   - sprints / iterations
%   - outils de planification (Trello, Jira, GitHub Projects...)
%   - planning previsionnel (diagramme de Gantt)

[\`A compl\'eter : m\'ethodologie adopt\'ee (ex. Scrum) et planning du
projet, \'eventuellement illustr\'e par un diagramme de Gantt.]

% \begin{figure}[H]
%     \centering
%     \includegraphics[width=0.9\textwidth]{images/gantt.png}
%     \caption{Diagramme de Gantt du projet}
%     \label{fig:gantt}
% \end{figure}

\section{Conclusion}
% Resumer en 2-3 lignes ce qui a ete vu (contexte, problematique, solution)
% et annoncer la transition vers le chapitre 2 (analyse des besoins).

[\`A compl\'eter]

\clearpage
